\documentclass[11pt,a4paper]{article}

\usepackage[margin=2.5cm]{geometry}
\usepackage{amsmath}
\usepackage{amssymb}
\usepackage{bm}
\usepackage{booktabs}
\usepackage{hyperref}

\title{Taichi MLS-MPM Dough Model}
\author{}
\date{}

\begin{document}
\maketitle

\section{State Variables}

The simulation represents dough with material particles coupled to a background grid.
For each particle $p$ the state is
\[
  \bm{x}_p,\quad \bm{v}_p,\quad \bm{F}_p,\quad \bm{C}_p,\quad J_{p,p}.
\]
Here $\bm{x}_p$ is position, $\bm{v}_p$ is velocity, $\bm{F}_p$ is the deformation
gradient, $\bm{C}_p$ is the APIC affine velocity matrix, and $J_{p,p}$ is the
accumulated plastic volume ratio. Grid nodes are indexed by $i$ and store
\[
  \bm{v}_i,\quad m_i .
\]

\section{Particle Initialization}

Particles are sampled inside an ellipsoid centered at
\[
  \bm{x}_0 = (0.5,\;0.32,\;0.5)^T
\]
with radii
\[
  \bm{r} = (0.14,\;0.055,\;0.12)^T .
\]
For a random point $\bm{q}$ inside the unit sphere,
\[
  \bm{x}_p = \bm{x}_0 + \bm{r}\odot \bm{q},
  \qquad
  \bm{v}_p = \bm{0},
  \qquad
  \bm{F}_p = \bm{I},
  \qquad
  \bm{C}_p = \bm{0},
  \qquad
  J_{p,p}=1 .
\]

\section{Material Parameters}

The Lam\'e parameters are computed from Young's modulus $E$ and Poisson ratio $\nu$:
\[
  \mu_0 = \frac{E}{2(1+\nu)},
  \qquad
  \lambda_0 = \frac{E\nu}{(1+\nu)(1-2\nu)} .
\]
If plastic volume hardening is enabled, stiffness is modified by
\[
  h_p = \exp\left(h_J(1 - J_{p,p})\right),
  \qquad
  \mu = h_p\mu_0,
  \qquad
  \lambda = h_p\lambda_0 ,
\]
where $h_J$ is the \texttt{jp-hardening} coefficient. With $h_J=0$, $J_{p,p}$
is tracked but does not change the stiffness.

\section{MLS-MPM Transfer}

Let the grid spacing be $\Delta x$ and $\Delta x^{-1} = n_\mathrm{grid}$.
For particle $p$,
\[
  \bm{g}_p = \bm{x}_p / \Delta x,
  \qquad
  \bm{b}_p = \lfloor \bm{g}_p - 0.5 \rfloor,
  \qquad
  \bm{f}_p = \bm{g}_p - \bm{b}_p .
\]
The one-dimensional quadratic B-spline weights are
\[
  w_0(f)=\frac{1}{2}(1.5-f)^2,
  \qquad
  w_1(f)=0.75-(f-1)^2,
  \qquad
  w_2(f)=\frac{1}{2}(f-0.5)^2 .
\]
For a neighboring node offset $\bm{o}=(a,b,c)$, $a,b,c\in\{0,1,2\}$,
\[
  W_{p\bm{o}} = w_a(f_x)w_b(f_y)w_c(f_z),
  \qquad
  \bm{d}_{p\bm{o}} = (\bm{o}-\bm{f}_p)\Delta x .
\]

\section{Deformation Gradient Update}

Before particle-to-grid transfer, the deformation gradient is updated by the
local velocity gradient:
\[
  \bm{F}_p^{*} =
  \left(\bm{I}+\Delta t\,\bm{C}_p\right)\bm{F}_p^n .
\]

\section{SVD Clamp Plasticity}

If plasticity is enabled, the model performs singular value clamping:
\[
  \bm{F}_p^{*} = \bm{U}\bm{\Sigma}\bm{V}^T,
  \qquad
  \tilde{\sigma}_k =
  \operatorname{clamp}(\sigma_k,\sigma_{\min},\sigma_{\max}) .
\]
The projected deformation gradient is
\[
  \bm{F}_p^{n+1} =
  \bm{U}\tilde{\bm{\Sigma}}\bm{V}^T .
\]
In the implementation,
\[
  \sigma_{\min}=\texttt{plastic-min},
  \qquad
  \sigma_{\max}=\texttt{plastic-max}.
\]
Values outside this interval are treated as plastic deformation, so only the
clamped part remains elastic.

\section{Plastic Volume Ratio}

If $J_p$ tracking is enabled, the accumulated plastic volume ratio is updated as
\[
  J_{p,p}^{n+1} =
  \operatorname{clamp}
  \left(
    J_{p,p}^{n}
    \frac{\det(\bm{F}_p^{*})}{\det(\bm{F}_p^{n+1})},
    J_{p,\min},
    J_{p,\max}
  \right).
\]
Here $J_{p,\min}$ and $J_{p,\max}$ correspond to \texttt{jp-min} and
\texttt{jp-max}.

\section{Constitutive Model}

The model uses corotated elastic stress plus a simple viscous term. Let
\[
  J = \det(\bm{F}_p),
  \qquad
  \bm{F}_p = \bm{R}\bm{S}
\]
be the polar decomposition. The elastic stress term used in the code is
\[
  \bm{P}_e =
  2\mu(\bm{F}_p-\bm{R})\bm{F}_p^T
  +
  \lambda J(J-1)\bm{I}.
\]
The viscous term is
\[
  \bm{P}_v =
  \eta(\bm{C}_p+\bm{C}_p^T),
\]
where $\eta$ is \texttt{viscosity}. The total stress-like quantity used for
MLS-MPM transfer is
\[
  \bm{P} = \bm{P}_e + \bm{P}_v .
\]

\section{Particle-to-Grid Transfer}

The affine momentum contribution is
\[
  \bm{A}_p =
  -\Delta t\,V_p\,4\,\Delta x^{-2}\bm{P}
  +
  m_p\bm{C}_p .
\]
Grid mass and momentum are accumulated as
\[
  m_i \mathrel{+}= W_{pi}m_p,
\]
\[
  m_i\bm{v}_i \mathrel{+}=
  W_{pi}
  \left(
    m_p\bm{v}_p + \bm{A}_p\bm{d}_{pi}
  \right).
\]
After accumulation,
\[
  \bm{v}_i \leftarrow \frac{\bm{v}_i}{m_i}.
\]

\section{Grid Update}

Gravity is applied on grid nodes:
\[
  \bm{v}_i \leftarrow \bm{v}_i + \Delta t(0,g,0)^T .
\]
The floor is a plane at $y=y_f$. If a grid node is below the floor and moving
downward,
\[
  v_{i,y}\leftarrow 0,
  \qquad
  v_{i,x}\leftarrow 0.7v_{i,x},
  \qquad
  v_{i,z}\leftarrow 0.7v_{i,z}.
\]

\section{Tool Box Contact}

Each tool is an oriented box with center $\bm{c}$, rotation $\bm{Q}$, visible
half-size
\[
  \bm{h} = (0.05,\;0.05,\;0.05)^T,
\]
and collision half-size
\[
  \bm{h}_c = \bm{h} + \epsilon_c,
\]
where $\epsilon_c$ is \texttt{tool-contact-padding}. A point $\bm{x}$ is
converted to tool-local coordinates by
\[
  \bm{q} = \bm{Q}^T(\bm{x}-\bm{c}).
\]
Collision is active when
\[
  |q_x| < h_{c,x},\qquad
  |q_y| < h_{c,y},\qquad
  |q_z| < h_{c,z}.
\]
The contact normal is chosen from the smallest penetration axis. For relative
velocity
\[
  \bm{u} = \bm{v}-\bm{v}_{tool},
\]
the inward normal component is removed:
\[
  \bm{u} \leftarrow
  \begin{cases}
    \bm{u} - (\bm{u}\cdot\bm{n})\bm{n},
      & \bm{u}\cdot\bm{n}<0,\\
    \bm{u},
      & \text{otherwise}.
  \end{cases}
\]
The remaining relative velocity is damped by friction and absorption:
\[
  \bm{v}\leftarrow
  \bm{v}_{tool}
  +
  \alpha_f(1-\alpha_a)\bm{u},
\]
where $\alpha_f$ is \texttt{tool-contact-friction} and $\alpha_a$ is
\texttt{tool-contact-absorption}. The same idea is applied on grid nodes and
again as a particle-level projection to reduce tunneling.

\section{Grid-to-Particle Transfer}

Particle velocity and affine velocity matrix are interpolated back from the
grid:
\[
  \bm{v}_p^{n+1} =
  \sum_i W_{pi}\bm{v}_i,
\]
\[
  \bm{C}_p^{n+1} =
  4\Delta x^{-1}
  \sum_i W_{pi}\bm{v}_i\bm{d}_{pi}^T .
\]
Global numerical damping is then applied:
\[
  \bm{v}_p^{n+1} \leftarrow
  \gamma_v \bm{v}_p^{n+1},
\]
where $\gamma_v$ is \texttt{velocity-damping}. Finally,
\[
  \bm{x}_p^{n+1} =
  \bm{x}_p^n + \Delta t\bm{v}_p^{n+1}.
\]

\section{Additional Plastic Damping}

If a particle yielded during SVD clamping, additional damping is applied:
\[
  \bm{v}_p \leftarrow \gamma_p\bm{v}_p,
  \qquad
  \bm{C}_p \leftarrow \gamma_C\bm{C}_p,
\]
where $\gamma_p$ is \texttt{plastic-velocity-damping} and $\gamma_C$ is
\texttt{plastic-affine-damping}. This is a pragmatic numerical term that reduces
rubber-like rebound after plastic deformation.

\section{Interpretation}

The current material should be described as a simplified MLS-MPM dough-like
model:
\[
  \text{corotated elasticity}
  +
  \text{viscous stress}
  +
  \text{SVD clamp plasticity}
  +
  \text{optional }J_p\text{ hardening}.
\]
It is inspired by common snow/plastic MPM formulations, but it is not yet a
calibrated food-dough rheology model.

\end{document}
